\chapter{Preface}

This manual presents the physical models and methods implemented in \pkg{}.
Each method is introduced through its Hamiltonian, governing equations,
approximations and physical interpretation. Code examples connect the equations
to the public Python interface. They are tested regularly, in the order in which
they appear; an example may therefore use variables defined earlier.

\paragraph{Three documents.}
This physics manual is one of three documents that are kept consistent with each other:
\begin{itemize}
\item the \emph{physics manual} explains what is computed and what it means;
\item the \emph{user and developer manual}, the Markdown pages in \code{docs/} (also built as a
website), explains how to install the code and how to run, check and save calculations; it lists
every option, result field and command-line argument, and describes how to extend the code;
\item the \emph{tutorials} in \code{examples/} are runnable calculations with checked results
that accompany the worked examples of \cref{part:examples}.
\end{itemize}
Each worked example names its tutorial, and the \enquote{In the code} boxes name the relevant
page of the user manual.

\paragraph{Who should read what.}
The opening chapters introduce superconducting impurities starting from basic
BCS theory and second quantisation. \Cref{ch:background} develops the
Andreev-reflection picture and the interpretation of quasiparticle occupations; \cref{ch:model} introduces the interacting
dot. \Cref{ch:correlations} explains YSR states, Kondo physics and screening, and
\cref{ch:experiments} connects them to spectroscopy and qubits. Short interpretation
exercises accompany the figures. \Cref{ch:static} compares solution methods before
describing the two numerical backends. Readers familiar with quantum-dot Josephson junctions
can start with the conventions and \cref{ch:practical}, which explains how to set up
a calculation, assess its accuracy and save the results before introducing periodic driving.
\Cref{ch:floquet} introduces the physics of periodically driven quantum systems
and develops Floquet theory from its basic principles.
\Cref{ch:driving,ch:drivenmethods} explain how drives enter the junction model and
how the code computes quasienergies; \cref{ch:probes} explains how a weak probe induces
transitions between Floquet states. \Cref{part:examples} then builds up the physics
step by step. An isolated dot, a driven two-level system, the high-frequency limit and a
chapter on numerical convergence come first. They are followed by the phase-driven and the
gate-driven junction, two-tone spectroscopy of the pair transition, photon-assisted
tunnelling, and electron spin resonance of the Zeeman-split doublet.
\Cref{ch:equations-to-code} maps the equations onto the code, and \cref{app:derivations} contains
the derivations that would interrupt the main text.

\paragraph{Conventions in brief.}
Energies are in units of the gap $\Delta$ and $\hbar=1$; $\Omega$ is then also an
energy. The bandwidth parameter is the normal-state \emph{half}-bandwidth $D$. The
hybridization $\Gamma$ is the \emph{total} hybridization summed over leads. Phases are in
radians, and the reference junction carries the phases $\mp\phi/2$ on its left and right
leads. Periodic Hamiltonians are expanded as $H(t)=\sum_k H_k\,\ee^{\ii k\Omega t}$.
\Cref{tab:conventions} on page~\pageref{tab:conventions} lists all conventions.

\paragraph{Boxes.}
Three kinds of boxes appear throughout:

\begin{physics}
These boxes explain the physical interpretation of a result, using equations where needed.
\end{physics}
\begin{incode}
These boxes identify the functions, options and result fields associated with a concept.
\end{incode}
\begin{caution}
These boxes highlight limitations and common numerical or interpretational errors.
\end{caution}

\paragraph{Reproducibility.}
Every figure in \cref{part:examples} is produced by a script in
\code{docs/manual/scripts}. The scripts store the numerical results, together with the
code version and numerical environment, in \code{docs/manual/data}; the figures are then
rendered from these records. \Cref{sec:reproduce} explains how to rebuild the figures and
this document.

\paragraph{Scope and status.}
The Floquet part of the code computes quasienergies, harmonic weights and Floquet modes. It
also computes the matrix elements of probe operators between static or Floquet states,
within one symmetry sector or between sectors, and the resulting spectral lines. Two
extensions are planned: the populations of the Floquet states in the presence of
dissipation, and a treatment of the continuum beyond the finite bath. The manual is
extended together with the code.
